\documentclass[11pt]{article}
\usepackage[a4paper,margin=1in]{geometry}
\usepackage{amsmath,amssymb}

\title{Subspace Alignment Training Summary}
\date{}

\begin{document}
\maketitle

\paragraph{Base objective.}
Training follows the standard SiT velocity objective. We sample
\[
\tau \sim \mathcal{U}(0,1), \qquad x_1 \sim \mathcal{N}(0,I),
\]
form
\[
x_\tau = (1-\tau)x_1 + \tau x_0,
\qquad
v^\star = x_0 - x_1,
\]
and optimize
\[
\mathcal{L}_{\mathrm{gen}} = \left\| v_\theta(x_\tau,\tau,y) - v^\star \right\|_2^2.
\]

\paragraph{Alignment target.}
The input latent is first patchified,
\[
x_0 \in \mathbb{R}^{N \times d_{\mathrm{in}}},
\]
then projected by the model patch embedding layer, before positional encoding:
\[
T_0 = \mathrm{PatchEmbed}(x_0) \in \mathbb{R}^{N \times D}.
\]
This is what we call \emph{layer 0}. It already lies in the hidden space of the model, but it has not yet absorbed absolute positional bias.

\paragraph{Aligned source layer.}
For a backbone layer $\ell$, let
\[
H_\ell \in \mathbb{R}^{N \times D}
\]
denote its token matrix. If \texttt{align-layer=random} is enabled, then $\ell$ is resampled every iteration. The target remains $T_0$, but the aligned source layer changes from step to step.

\paragraph{Centering.}
Before extracting any subspace, both representations are centered over the patch axis:
\[
\bar{H}_\ell = H_\ell - \frac{1}{N}\mathbf{1}\mathbf{1}^\top H_\ell,
\qquad
\bar{T}_0 = \mathrm{sg}\!\left(T_0 - \frac{1}{N}\mathbf{1}\mathbf{1}^\top T_0\right).
\]
This removes the global token bias and keeps only relative spatial variation across patches.

\paragraph{TSVD subspace extraction.}
The TSVD branch extracts a $k$-dimensional subspace over patches. In principle one may compute
\[
\bar{H}_\ell = U_H \Sigma_H V_H^\top
\]
and keep the first $k$ columns of $U_H$. The exact reference computation uses the equivalent Gram formulation
\[
G_H = \bar{H}_\ell \bar{H}_\ell^\top = U_H \Sigma_H^2 U_H^\top,
\]
and similarly
\[
G_T = \bar{T}_0 \bar{T}_0^\top = U_T \Sigma_T^2 U_T^\top.
\]
Therefore, the top eigenvectors of $G_H$ and $G_T$ are exactly the left singular subspaces of $\bar{H}_\ell$ and $\bar{T}_0$. We set
\[
Q_H = \mathrm{TopEig}_k(G_H), \qquad Q_T = \mathrm{TopEig}_k(G_T),
\]
with $Q_H,Q_T \in \mathbb{R}^{N \times k}$ and orthonormal columns.

\paragraph{Why this is equivalent to full TSVD.}
If one ran full TSVD on $X \in \mathbb{R}^{N \times D}$ and wrote
\[
X = U \Sigma V^\top,
\]
then the left singular vectors used by TSVD are exactly the eigenvectors of
\[
X X^\top = U \Sigma^2 U^\top.
\]
So, for the purpose of extracting the top-$k$ patch subspace, \emph{full SVD on $X$} and \emph{eigendecomposition of $X X^\top$} return the same object. In the generic case they recover the same ordered directions. When singular values are repeated, the basis inside that repeated eigenspace is not unique, but the span is still the same. Since the loss only depends on the projector
\[
P = Q Q^\top,
\]
the optimization target is unchanged even if the basis vectors themselves rotate inside the same subspace.

To reduce the chance of spectral degeneracy in the train-side TSVD, we add a very small Gaussian perturbation before extracting $Q_H$:
\[
\tilde{H}_\ell = \bar{H}_\ell + \alpha \,\mathrm{RMS}(\bar{H}_\ell)\,\varepsilon,
\qquad
\varepsilon \sim \mathcal{N}(0,I).
\]
This perturbation is only used on the train side. The target branch stays clean and is stop-gradient.

For exact TSVD training, the forward pass uses the exact Gram eigenspace above. To avoid undefined eigenvector gradients at repeated or nearly repeated eigenvalues, the backward pass is routed through a stabilized subspace-iteration surrogate. Starting from a fixed orthonormal basis $Q^{(0)}$, the surrogate repeatedly applies the patch Gram operator without materializing the full Gram matrix:
\[
Z^{(r)} = X^\top Q^{(r)}, \qquad
\tilde{Q}^{(r+1)} = X Z^{(r)}, \qquad
Q^{(r+1)} = \mathrm{QR}(\tilde{Q}^{(r+1)}).
\]
The default train solver is therefore exact in the forward objective, while the surrogate only controls the gradient path. A faster purely iterative train solver remains available as an optional speed/debug setting, and validation keeps the exact Gram eigensolver by default.

\paragraph{TSVD warmup.}
When TSVD alignment is enabled, the first
\[
K_{\mathrm{warm}} = 25{,}000
\]
training steps are run with the base SiT objective only. During this warmup, the TSVD branch is skipped entirely:
\[
\lambda_{\mathrm{align}}^{\mathrm{eff}}(s) = 0,
\qquad s < K_{\mathrm{warm}}.
\]
After warmup, the configured alignment weight is applied:
\[
\lambda_{\mathrm{align}}^{\mathrm{eff}}(s) = \lambda_{\mathrm{align}},
\qquad s \ge K_{\mathrm{warm}}.
\]
This lets the backbone first move away from the highly symmetric initialization regime before TSVD eigenvectors are used in the auxiliary loss.

\paragraph{Projector loss.}
For either TSVD or QBA, the aligned object is the patch subspace, not the basis vectors themselves. The canonical object is the projector
\[
P = QQ^\top.
\]
The raw spatial discrepancy is
\[
\mathcal{L}_{\mathrm{spatial}}
=
\left\| P_H - P_T \right\|_F^2.
\]
Because $Q_H$ and $Q_T$ are orthonormal, this can be evaluated more cheaply as
\[
\mathcal{L}_{\mathrm{spatial}}
=
2k - 2\,\mathrm{tr}(P_H P_T)
=
2k - 2\left\|Q_H^\top Q_T\right\|_F^2.
\]
This expression is mathematically identical to the projector Frobenius loss, but avoids explicitly materializing the full $N \times N$ projectors.

\paragraph{QBA.}
The QBA branch replaces SVD/eigendecomposition on both sides by symmetric
learnable bottlenecks:
\[
W_H,W_T \in \mathbb{R}^{D \times k},
\qquad
Z_H = \bar{H}_\ell W_H,
\qquad
Z_T = \bar{T}_0 W_T.
\]
The layer-0 tokens remain stop-gradient through $\bar{T}_0$, but the target
branch has its own learned bottleneck $W_T$ rather than a fixed basis. Thus both
branches use the same project-then-orthonormalize mechanism before projector
matching.

Both bottleneck representations are then orthonormalized by thin QR,
\[
Z_H = Q_H R_H,
\qquad
Z_T = Q_T R_T,
\]
and the same spatial loss above is applied to the resulting subspaces. The QR
step is the bottleneck whitening/orthonormalization used by QBA: the $R$
factors are discarded, while $Q_H$ and $Q_T$ provide scale- and
correlation-normalized bases for projector matching. We regularize both learned
bottlenecks:
\[
\mathcal{L}_{\mathrm{ortho}} =
\left\| W_H^\top W_H - I \right\|_F^2
+
\left\| W_T^\top W_T - I \right\|_F^2.
\]
Hence
\[
\mathcal{L}_{\mathrm{QBA}} =
\mathcal{L}_{\mathrm{spatial}}
+
\lambda_{\mathrm{ortho}}\mathcal{L}_{\mathrm{ortho}}.
\]

\paragraph{Final objective.}
The full training loss is
\[
\mathcal{L}
=
\mathcal{L}_{\mathrm{gen}}
+
\lambda_{\mathrm{align}}\mathcal{L}_{\mathrm{align}},
\]
where $\mathcal{L}_{\mathrm{align}}$ is either the TSVD projector loss or the QBA loss above. In words, the model is still trained to solve the diffusion task, but it is additionally encouraged to preserve the same coarse patch geometry as layer 0.

\end{document}
